\chapter{Machine-learning methods in depth}
\label{ch:ml-appendix}

\begin{quote}\itshape
This appendix explains \emph{how} each learning method in the document actually works, one
level below the equations of \cref{sec:ai-math,sec:stats-math}. It assumes only the
six-idea primer of \cref{sec:ml-primer}; a reader who wants the mechanics --- what a network
computes, how it is trained, why each choice was made and where it breaks --- will find them
here. Nothing new is claimed; this is the reference layer the main text points to.
\end{quote}

\section{How a neural network is trained}
\label{sec:ml-nn}

Every neural model here (the GRU, the autoencoder, the surrogate) is trained by the same
loop, worth stating once in full.

\paragraph{The forward pass.} A network is a function $f_\theta(\mathbf{x})$ built by
alternating linear maps and simple nonlinearities. A single layer computes
$\mathbf{a} = \sigma(\mathbf{W}\mathbf{x} + \mathbf{b})$: multiply the input by a weight
matrix $\mathbf{W}$, add a bias $\mathbf{b}$, and apply an element-wise nonlinearity
$\sigma$ (here GELU, a smooth relative of the rectifier). Stacking $L$ such layers lets the
network represent very general input--output relationships; the adjustable numbers
$\theta = \{\mathbf{W}^{(1)}, \mathbf{b}^{(1)}, \ldots\}$ are the parameters.

\paragraph{The loss.} Training needs a scalar measure of wrongness. For regression (RUL,
the surrogate) it is the mean squared error $\mathcal{L} = \frac1N\sum_i \|f_\theta(\mathbf{x}_i)
- y_i\|^2$; for classification (diagnosis) it is the cross-entropy
$\mathcal{L} = -\frac1N\sum_i \log f_\theta(\mathbf{x}_i)_{y_i}$, which rewards putting
probability mass on the true class. Class weighting multiplies each term by the inverse
class frequency so rare fault classes are not ignored.

\paragraph{Backpropagation.} To reduce $\mathcal{L}$ we need its gradient with respect to
every parameter, $\partial\mathcal{L}/\partial\theta$. Backpropagation is the chain rule
applied layer by layer from the output backwards: the error at the output is propagated
back through each layer's local derivative, accumulating the contribution of every weight.
It costs one backward pass per forward pass --- linear, not exponential, in the number of
parameters, which is what makes deep networks trainable at all.

\paragraph{Gradient descent and Adam.} Each step nudges the parameters against the gradient,
$\theta \leftarrow \theta - \eta\, \partial\mathcal{L}/\partial\theta$, with learning rate
$\eta$. Plain descent is slow and sensitive to $\eta$; the \emph{Adam} optimizer used here
keeps running averages of the gradient (momentum) and of its square (per-parameter scaling),
so each weight gets an adaptive step size. The gradient is computed not on the whole dataset
but on random \emph{mini-batches}, which is faster and adds a helpful noise; one sweep
through the data is an \emph{epoch}, and training runs for tens of them.

\paragraph{Why it does not simply memorize.} A network with enough parameters can fit its
training set exactly and generalize terribly (\cref{sec:ml-primer}). The defenses used here
are the honest split (all reported numbers come from engines never trained on), early-ish
stopping (a bounded epoch budget), and keeping the models small --- the RUL GRU has on the
order of $10^4$ parameters, not $10^7$. On this dataset bigger networks did not help, which
is itself informative: the signal, not model capacity, is the ceiling.

\section{The GRU, mechanism and motivation}
\label{sec:ml-gru}

The prognosis and F7 isolation models are gated recurrent units
(\cref{eq:gru}). A recurrent network reads a sequence one step at a time, carrying a hidden
state $\mathbf{h}_n$ that summarizes everything seen so far --- the natural fit for an
engine's deviation history.

\paragraph{The problem it solves.} A plain recurrent network multiplies the hidden state by
the same matrix at every step; over hundreds of steps that repeated multiplication makes
gradients either vanish or explode, so the network cannot learn long-range dependence (the
early-life trend that matters for late-life RUL). The GRU fixes this with \emph{gates}:
learned, data-dependent valves that decide how much of the state to keep versus rewrite.

\paragraph{The gates, in words.} At each step (\cref{eq:gru}): the \emph{update gate}
$\mathbf{z}_n$ interpolates between keeping the old state and adopting a freshly computed
candidate --- near zero it preserves memory across a quiet stretch, near one it rewrites on
new evidence. The \emph{reset gate} $\mathbf{r}_n$ decides how much of the past to consult
when forming that candidate. Because the update gate can hold the state unchanged, the
gradient can flow across many steps without vanishing --- the network can ``remember'' the
early trend. The whole advantage over the classical Holt smoother
(\cref{eq:holt}) is that the GRU's forgetting rate is \emph{learned, vector-valued and
context-dependent} rather than a single fixed constant.

\paragraph{Training a recurrent network.} Backpropagation through a sequence is
``backpropagation through time'': the network is unrolled over the window (here 32--96
steps), and the chain rule runs back through every step. Sequences are cut into fixed-length
windows and down-sampled (every 10th--30th cycle) both to bound the unroll length and because
adjacent flights are near-duplicates.

\section{The autoencoder and why it failed}
\label{sec:ml-ae}

The v0 detector was an autoencoder (\cref{eq:autoencoder}): an encoder compresses a window to
a low-dimensional code, a decoder reconstructs it, and training minimizes reconstruction
error on \emph{healthy} data only. The premise is that a model fluent only in healthy
patterns will reconstruct a faulty window poorly, so the reconstruction error is an anomaly
score. It failed for a reason visible in its own objective: with a code wide enough to be
expressive, the cheapest way to minimize reconstruction error is to learn an approximate
\emph{identity map} --- copy the input through --- which reconstructs faulty and healthy
windows equally well and flattens the score. The lesson generalizes: reconstruction-based
anomaly detection needs an information bottleneck tight enough to forbid the identity
shortcut, and the winning Mahalanobis detector avoids the trap entirely by scoring against
the healthy distribution directly rather than learning to reconstruct.

\section{Mahalanobis detection with covariance shrinkage}
\label{sec:ml-maha}

The detector that works (\cref{eq:mahalanobis}) asks how atypical a feature vector is
relative to healthy behaviour, measured in ``whitened'' units that undo the channels'
different scales and correlations. The one subtlety is the covariance estimate. With more
features than comfortable relative to the sample size, the empirical covariance
$\mathbf{S}$ is nearly singular: its smallest eigenvalues are underestimated, and inverting
it magnifies noise along those directions, manufacturing false alarms. Ledoit--Wolf
shrinkage (\cref{eq:ledoitwolf}) blends $\mathbf{S}$ toward a scaled identity by an amount
$\rho$ computed in closed form to minimize expected error --- no tuning knob. The result is a
well-conditioned covariance and a stable distance. Detection then requires several consecutive
exceedances of a healthy-data percentile, turning a per-snapshot score into an alarm with a
controlled false-positive rate.

\section{Gradient-boosted trees for diagnosis}
\label{sec:ml-gbdt}

The isolation classifier is a histogram gradient-boosting machine (\cref{eq:gbdt}): an
ensemble of shallow decision trees built one at a time, each new tree fitted to the gradient
of the loss left by the ensemble so far --- so each tree corrects the current mistakes.
``Histogram'' means continuous features are bucketed into a few dozen bins before split
search, which makes training fast. Boosted trees remain the strongest off-the-shelf learner
on small \emph{tabular} problems (hundreds of engineered step features, not raw sequences),
which is why diagnosis got trees while the sequence tasks got the GRU. Their limit here is
not the model but the data: on the confusable subset the information is simply absent
(\cref{sec:res-h2}), and no tree depth recovers it.

\section{Conformal prediction, in detail}
\label{sec:ml-conformal}

Split conformal (\cref{eq:conformal}) turns any point predictor into an interval with a
finite-sample coverage guarantee, assuming only that calibration and test points are
\emph{exchangeable} (their joint distribution is order-independent --- weaker than
independence, and satisfied by an engine-level split). The recipe: on a held-out calibration
set collect the absolute residuals, take their $\lceil(1-\alpha)(n_c+1)\rceil$-th smallest
value $\hat{q}$, and predict $\hat{y}(x) \pm \hat{q}$. The guarantee
$\Pr[y \in \hat{y}\pm\hat{q}] \ge 1-\alpha$ follows because the new point's residual is
equally likely to fall in any rank among the calibration residuals; the $(n_c+1)$ counts the
new point joining the ranking, which is what makes the bound exact rather than asymptotic.

\paragraph{Prediction sets for classification.} The isolation variant (adaptive prediction
sets, \cref{ch:tomography}) applies the same logic to a classifier: sort the predicted class
probabilities, accumulate them until the true class is included, and calibrate the
accumulation threshold on validation data. The output is a \emph{set} of classes whose size
grows when the model is unsure --- ambiguity becomes an explicit, guaranteed-coverage object
rather than a silent error, which is exactly what the F7 isolation sets and the F10
certificate (\cref{ch:breakthrough}) exploit.

\section{SHAP attributions}
\label{sec:ml-shap}

To ask \emph{why} a classifier made a call, SHAP assigns each input feature a Shapley value
(\cref{eq:shapley}): its average contribution to the prediction over all orders in which
features could be revealed, with the rest marginalized over background data. Shapley values
are the unique attribution satisfying three fairness axioms from cooperative game theory
(the contributions sum to the prediction minus a baseline; identical features get equal
credit; a feature that never changes the output gets zero). The exact computation is
exponential in the number of features, so a permutation sampler estimates it. In this
project the attribution vector is not the end product but the input to the
Physics-Consistency Score (\cref{sec:pcs}), which projects it through the influence matrix
and checks whether the model's stated reasoning points at the physically correct component.

\section{Tuning: the tree-structured Parzen estimator}
\label{sec:ml-tpe}

Both EHM families were tuned by Optuna's tree-structured Parzen estimator (TPE,
\cref{sec:stats-math}). After a few random trials it splits the tried configurations into a
good group and the rest, fits a probability density over each, and proposes the next
configuration where the ratio (good density)/(rest density) is largest --- concentrating
trials where good settings cluster while still exploring. On a 14--16-dimensional mixed
search space this finds strong configurations in $\sim$50 trials, far fewer than grid search
would need, and --- critically for fairness --- both families receive the same budget.

\section{The neural surrogate: residual learning}
\label{sec:ml-surrogate}

The differentiable twin (\cref{sec:f8-l1}) is a neural network that predicts the engine's
channel values from its health state and operating condition. The architectural choice that
made it pass was to predict not the value but the \emph{residual} above the F1 linearization:
$\text{surrogate} = \text{linear model} + \text{MLP correction}$. This inherits the linear
model's near-perfect accuracy where physics is nearly linear (so the network need not
re-learn it) and spends the network's capacity only on the curvature the linear model
misses. A pure network, asked to learn the whole map, could not match the linearization's
precision on the smoothest channel (N2); the residual form fixed it at once --- the same
``keep the physics, learn the correction'' pattern that, applied to features rather than
outputs, \emph{failed} for the RUL hybrid (H4). The difference is instructive: residual
learning helps when the physics is a good first approximation of the target (channel values),
and hurts when the physics-derived quantity is a noisy, redundant feature of an already-solved
problem (RUL).

\section{Learned MOPA: physics inversion, then learned fusion}
\label{sec:ml-mopa}

The F7 isolation model (\cref{sec:tomo-learner}) is the clearest illustration of where
learning belongs relative to physics. A GRU fed raw deviation sequences failed (12\,\%
isolation): asked to learn the entire inverse physics from a few hundred windows, it had too
little data. A GRU fed the \emph{per-flight regularized WLS projections} --- the physics
inversion done first, analytically, at each flight's condition --- then learned only to
\emph{fuse} that sequence of noisy estimates over time, and tripled classical stacking.
The division of labour is the lesson: let the known physics do the inversion it does well,
and let the network do the temporal fusion it does well; neither the pure-classical nor the
pure-neural extreme wins.
